\section{理论设定}

本节简要介绍我们想要重正化的准 PDF 的核子矩阵元，并解释三类 PDF——非极化、螺旋度和横向性——
的重正化处方。核子矩阵元计算的细节见文献 Refs.~\cite{Alexandrou:2015rja,Alexandrou:2016jqi}。

%#############################################################
\vspace*{0.25cm}
\subsection{核子矩阵元}
\vspace*{0.15cm}
%#############################################################

我们考虑含一条直 Wilson 线的非局域费米子算符的矩阵元，记为 $h_\Gamma(P_3,z)$。变量 $z$ 是 Wilson 线的长度，
$P_3$ 是核子动量，取在与 Wilson 线相同的方向上。准 PDF 可由下列局域矩阵元的傅里叶变换计算：
\begin{equation}
\label{MElements}
h_\Gamma(P_3,z)=\langle N |\bar{\psi}(0,z)\,\Gamma \,W_3(z) \psi(0,0) |N\rangle,
\end{equation}
其中 $|N\rangle$ 是空间动量 $\vec{P}{=}(0,0,P_3)$ 沿 $3$ 方向的核子态，$W_3(z)$ 是同方向、长度为 $z$ 的
Wilson 线。$\Gamma$ 表示算符插入的 Dirac 结构：$\gamma_\mu$（非极化）、$\gamma_\mu\cdot\gamma_5$（螺旋度）、
$\sigma_{\mu \nu}$（横向性）。在文献已有的工作中，$\gamma_\mu$ 取在沿 Wilson 线的方向。原则上也可以选择
$\gamma_\mu$ 垂直于 Wilson 线的方向；此时非极化算符没有混合，而螺旋度和横向性会混合。例如，如文献
~\cite{Radyushkin:2016hsy} 所讨论的，把 $\gamma$ 矩阵取在时间方向对更快地收敛到物理 PDF 很重要。
我们最近的工作~\cite{Constantinou:2017sej} 表明，$\vec P_3 {\perp} \vec z$ \,\,（$\vec P_3 {\parallel} \vec z$）
对非极化（螺旋度和横向性）情形是最理想的。

\medskip
为计算裸矩阵元，我们采用文献~\cite{Alexandrou:2016jqi} 的设定。我们考虑由 ETMC 产生的
一个动力学 $N_f{=}2{+}1{+}1$ 扭转质量费米子系综~\cite{Baron:2010bv}，体积 $32^3 {\times} 64$，
格距 $a {\approx}0.082$\,fm~\cite{Carrasco:2014cwa}，裸扭转质量 $a\mu {=} 0.0055$，对应约 375 MeV 的
$\pi$ 介子质量。我们在 1000 个规范组态上进行计算，包含 15 个前向传播子和 2 个随机传播子，
即总共 30000 次测量。我们将给出动量 $P_3=\frac{6\pi}{L}$（物理单位下约 1.4 GeV）的结果。
在该矩阵元计算中，对核子内插场采用了高斯涂抹（Gaussian smearing）~\cite{Alexandrou:2016jqi}。

%#############################################################
\vspace*{0.25cm}
\subsection{重正化方案}
\label{sec:scheme}
\vspace*{0.15cm}
%#############################################################

这里我们讨论一个完全非微扰的重正化处方，它将消除准 PDF 矩阵元中继承的全部发散以及混合——后者如文献
~\cite{Constantinou:2017sej} 的微扰分析所指出。简言之，所提出的重正化程序：
\begin{itemize}
\item[\bf 1] 消除可重新求和为一个乘性指数因子 $e^{-\delta m |z|/a + c |z|}$ 的线性发散。系数 $\delta m$
代表发散的强度，预期与算符无关，因为它只与 Wilson 线有关。$c$ 是一个任意标度~\cite{Sommer:2015hea}，
可通过这样的重正化处方加以固定；
\item[\bf 2] 消去相对调节子的对数发散 $\log(a\bar\mu_0)$，其中 $\bar\mu_0$ 是 RI$'$ 重正化标度；
\item[\bf 3] 施行与格点正规化相联系的必要的有限重正化；
\item[\bf 4] 消除非极化算符中出现的混合，因为其裸矩阵元是非极化准 PDF 与 twist-3 标量算符的线性组合。
二者可以通过构造 $2{\times}2$ 混合矩阵来分离。
\end{itemize}

\medskip
我们采用可以非微扰应用的重正化方案，即 RI$'$~方案 \cite{Martinelli:1994ty}。我们在动量空间的设定下、
在外部夸克态之间计算所研究算符的顶点函数，算符定义为：
\be
{\cal O}_\Gamma = \overline \psi(x) \, \Gamma \,\mathcal{P}\, e^{i\,g\,\int_{0}^{z} A(\zeta) d\zeta}\, \psi(x+z\hat{\mu})\,,
\ee
其中 $\Gamma=\gamma_\mu,\,\gamma_\mu\cdot\gamma_5,\,\sigma_{\mu\nu}$（$\nu\neq\mu$）。上式中出现的指数的路径排序
在格点上成为一系列按路径排序的规范连接（gauge link）。重正化函数（$Z$ 因子）依赖 Wilson 线的长度，
因此我们对每个 $z$ 值分别进行计算。通常 $z$ 可以达到格点空间延展的一半。

\medskip
重正化处方遵循为局域算符发展起来的程序的思路，顶点函数的构造在文献~\cite{Alexandrou:2010me} 中描述。
局域算符与 Wilson 线算符重正化的区别在于后者出现的线性发散。然而，无需把这个发散从乘性重正化中分离出来，
因此下面描述的技术可以同时成功地提取两种贡献。

%%%%%%%%%%%%%%%%%%%%%%%%%%%%%%%%%%%%%%%%%%%%%
\vspace*{0.55cm}
%\newpage
\centerline{\bf\underline{\textit{螺旋度与横向性准 PDF}}}
\vskip 0.25cm
%%%%%%%%%%%%%%%%%%%%%%%%%%%%%%%%%%%%%%%%%%%%%

我们首先给出在没有任何混合时一般 Wilson 线算符的方法。这适用于螺旋度与横向性准 PDF，
只要它们的 Dirac 结构选在沿 Wilson 线的方向。Wilson 线算符的重正化函数 $Z_{\cal O}$
通过施加如下条件来提取：
\be
Z_q^{-1}\,Z_{\cal O}(z)\,\frac{1}{12} {\rm Tr} \left[{\cal V}(p,z) \left({\cal V}^{\rm Born}(p,z)\right)^{-1}\right] \Bigr|_{p^2=\bar\mu_0^2} = 1\, ,
\label{renorm}
\ee
其中 $Z_q$ 是通过下式得到的夸克场的重正化函数：
\be
Z_q = \frac{1}{12} {\rm Tr} \left[(S(p))^{-1}\, S^{\rm Born}(p)\right] \Bigr|_{p^2=\bar\mu_0^2}  \,.
\label{Zq_cond}
\ee
迹对自旋与颜色指标求取，进入顶点函数的动量 $p$ 取为 RI$'$ 重正化标度 $\bar\mu_0$。
在式 (\ref{renorm}) 中，${\cal V}(p,z)$ 是算符的缩并顶点函数（amputated vertex function），
${\cal V}^{{\rm Born}}$ 是它的树图值，即对螺旋度算符 ${\cal V}^{{\rm Born}}(p,z) {=} i\gamma_3\gamma_5 \,e^{i p z}$。
同样，$S(p)$ 是费米子传播子，$S^{{\rm Born}}(p)$ 是其树图值。RI$'$ 标度 $\bar\mu_0$ 选择为其 $z$ 分量
与核子动量相同。这样的选择起到压制离散化效应的作用：不同类别的空间分量具有不同的离散化效应，
而形如 $(n_t,3,3,3)$ 的标度离散化效应较小~\cite{Alexandrou:2015sea}。我们检验了相对 Wilson 线
（空间方向）的对角（民主）动量与平行动量两种选择，即
$a\vec{\bar{\mu}}_0{=} \frac{2\pi}{L}\,(P_3,P_3,P_3)$ 与 $a\vec{\bar{\mu}}_0{=} \frac{2\pi}{L}\,(0,0,P_3)$，
分别称之为``对角''（diagonal）和``平行''（parallel）。后者预期有较大的格点伪影，因为比值
\be
\label{Phat}
\hat{P}{\equiv}\frac{\sum_\rho \bar\mu_{0_\rho}^4}{\left(\sum_\rho \bar\mu_{0_\rho}^2\right)^2}
\ee
高于对角动量的相应值。使用使该比值较小的重正化标度，对局域费米子算符已被证明是成功的
~\cite{Constantinou:2010gr,Alexandrou:2010me}。

\medskip
顶点函数 ${\cal V}(p)$ 含有与核子矩阵元相同的线性发散。这一点至关重要，因为它允许通过式 (\ref{renorm})
的重正化条件把指数连同乘性 $Z$ 因子一起提取出来。关于准 PDF 的可重正化性，
文献 Refs.~\cite{Ishikawa:2016znu,Ishikawa:2017faj} 证明它在微扰论各阶都是乘性的。
作者们考虑坐标空间中定义的准 PDF 并证明其乘性可重正化性，同样的结论对式 (\ref{MElements}) 的定义也成立。
这是因为前者额外要求一个一致的减除方案，以除去在极限 $\xi_z\rightarrow 0$ 下发散的所有项。
此外，双局域复合算符的重正化也在文献 Refs.~\cite{Ishikawa:2016znu,Ishikawa:2017faj} 中被研究。
基于以上论述，$Z_{\cal O}$ 可以因子化为
\be
Z_{\cal O}(z) = \overline{\cal Z}_{\cal O} \,\, e^{+\delta m |z|/a - c |z|}\,,
\ee
其中 $\overline{\cal Z}_{\cal O}$ 是算符的乘性 $Z$ 因子，$\delta m$ 是线性发散的强度。
该指数包含一项带任意标度 $c$ 的贡献，它可以取 $c|z|$ 的形式~\cite{Sommer:2015hea}。注意，
指数的符号与核子矩阵元中的相比是不同的，因为 $Z_{\cal O}$ 与顶点函数的逆相联系：
\be
Z_{\cal O} = \frac{Z_q}{\frac{1}{12} {\rm Tr} \left[{\cal V}(p) \left({\cal V}^{\rm Born}(p)\right)^{-1}\right] \Bigr|_{p=\bar\mu} }\,.
\ee
只要在同一组态系综上计算，$Z$ 因子的这种构造就说明了为什么用 $Z_{\cal O}$ 相乘能够成功消除
核子矩阵元中的幂发散。于是，上述重正化条件同时处理了所考虑矩阵元中存在的全部发散。
注意在没有 Wilson 线时（$z{=}0$），这些 $Z$ 因子退化为局域流的 $Z$ 因子，后者不含任何幂发散。
例如，对螺旋度算符，$Z_A(z{=}0)$ 就是轴向流的标准 $Z$ 因子。

\medskip
我们想指出，知道系数 $\delta m$ 能提供对幂发散强弱的洞见。可以通过静态势~\cite{Ishikawa:2016znu}
或文献~\cite{Constantinou:2017sej} 提出的技术追求这一方向。若提取出了线性发散，便可以应用
文献~\cite{Chen:2016fxx} 的匹配（其中包含系数 $\delta m$）。然而，仍有必要计算乘性 $Z$ 因子，
以治愈任何对数发散、施行有限重正化，以及固定任意标度 $c$。


%%%%%%%%%%%%%%%%%%%%%%%%%%%%%%%%%%%%%%%%%%%%%
\vspace*{0.75cm}
\centerline{\bf\underline{\textit{非极化准 PDF}}}
\vskip 0.25cm
%%%%%%%%%%%%%%%%%%%%%%%%%%%%%%%%%%%%%%%%%%%%%

如果 Dirac 结构取在与 Wilson 线相同的方向上，非极化准 PDF 的情形就需要特殊处理。
如文献~\cite{Constantinou:2017sej} 所证明的，这种选择下存在与 twist-3 标量算符的混合
\footnote{对扭转质量费米子，混合发生在矢量流与赝标流之间。}，必须予以消除。
由于手征对称性的破缺，这种混合出现在某些格点正规化中（特别是 Wilson 与扭转质量费米子），并且发现它是有限的。
我们建立如下记号：下标 $S$ 和 $V$ 分别表示标量算符与矢量（非极化）算符。相应的算符是：
\bea
{\cal O}_S \hspace*{-0.15cm}&=&\hspace*{-0.15cm} \overline \psi(x) \, \hat{1} \,\mathcal{P}\, e^{i\,g\,\int_{0}^{z} A(\zeta) d\zeta}\, \psi(x+z\hat{\mu})\,, \\[2ex]
{\cal O}_V \hspace*{-0.15cm}&=&\hspace*{-0.15cm} \overline \psi(x) \, \gamma_\mu \,\mathcal{P}\, e^{i\,g\,\int_{0}^{z} A(\zeta) d\zeta}\, \psi(x+z\hat{\mu})\,,
\eea
我们把它们的核子矩阵元记为 $h_S(P_3,z)$ 和 $h_V(P_3,z)$。
为了把这两个贡献从各自的裸矩阵元中分离出来，必须从下面的 $2{\times}2$ 矩阵计算乘性重正化和混合系数：
\begin{equation}
  \binom{{\cal O}^R_V(P_3,z)}{{\cal O}^R_S(P_3,z)} = \hat{Z}(z)\cdot
  \binom{{\cal O}_V(P_3,z)}{{\cal O}_S(P_3,z)}\,,
\end{equation}
其中 $\hat{Z}$ 是标量算符与矢量算符之间混合的矩阵：
\begin{equation}
 \hat{Z}(z) = \begin{pmatrix} Z_{VV}(z) & Z_{VS}(z) \\ Z_{SV}(z) & Z_{SS}(z) \end{pmatrix} \,.
\end{equation}
%
根据上述混合关系，重正化的非极化准 PDF $h^R_V(P_3,z)$ 通过下式与裸标量和裸非极化矩阵元相联系：
%
\begin{equation}
h^R_V(P_3,z) = Z_{VV}(z) \,\, h_V(P_3,z) + Z_{VS}(z) \,\, h_S(P_3,z) \,,
\label{h_R}
\end{equation}
%
其中 $Z_{VV}$ 和 $Z_{VS}$ 在某个确定的方案中计算。在本工作中，我们先采用 RI$'$ 方案，
然后在能量标度 $\bar\mu$ 下转换到 $\MSb$ 方案。$Z_{ii}$ 因子可以按照类似式 (\ref{renorm}) 的处方计算：
\be
\label{renormMIX}
Z_q^{-1}\,\hat{Z}(z)\, {\cal \hat{V}}(p,z)\Bigr|_{p=\bar\mu} = \hat{1}\,,
\ee
%
其中顶点函数矩阵 ${\cal \hat{V}}$ 的元素由迹给出：
\be
\left({\cal \hat{V}}(z) \right)_{ij}=  \frac{1}{12} {\rm Tr} \left[{\cal V}_i(p,z) \left({\cal V}_j^{\rm Born}(p,z)\right)^{-1}\right]\,,\quad i,j=S,V\,.
\ee
在上面的方程中，${\cal V}_i^{\rm Born}$ 是算符 ${\cal O}_i$ 的树图表达式。
这样，$\hat{Z}$ 的所有矩阵元都可以通过一组线性方程提取，这组方程可写成如下矩阵形式：
\be
Z_q^{-1}\,\begin{pmatrix} Z_{VV}(z) & Z_{VS}(z) \\[2ex] Z_{SV}(z) & Z_{SS}(z) \end{pmatrix}  \cdot
\begin{pmatrix}  \left({\cal \hat{V}}(z) \right)_{VV} & \left({\cal \hat{V}}(z) \right)_{SV} \\[2ex]
 \left({\cal \hat{V}}(z) \right)_{VS} &\left({\cal \hat{V}}(z) \right)_{SS} \end{pmatrix} =
\begin{pmatrix} 1& 0 \\[2ex] 0 & 1\end{pmatrix} \,.
\ee
从式 (\ref{h_R}) 可以看到，非极化准 PDF 重正化的一个主要组成部分是知道标量核子矩阵元 $h_S(P_3,z)$。
迄今为止还没有关于 $h_S(P_3,z)$ 的格点计算，因为在文献~\cite{Constantinou:2017sej} 的工作之前并未预料到会有这种混合。
因此，非极化准 PDF 的恰当重正化仍悬而未决。然而，当费米子作用量中含有 clover 项时，混合似乎受到极大压制。
这一点非常重要，因为我们目前正在物理 $\pi$ 介子质量的扭转质量 clover 改进费米子系综
~\cite{Abdel-Rehim:2015pwa,Abdel-Rehim:2015owa} 上计算准 PDF。


%#############################################################
\vspace*{0.25cm}
\subsection{转换到 $\MSb$ 方案}
\vspace*{0.15cm}
%#############################################################

含有限长 Wilson 线的费米子算符是方案和标度相关的。因此，当在 RI$'$ 标度 $\bar\mu_0$ 处
按式 (\ref{renorm}) 和 (\ref{renormMIX}) 得到 $Z$ 因子之后，必须把它们转换到标度为 $\bar\mu$ 的
$\MSb$ 方案。这一转换因子已在单圈连续微扰论中于文献~\cite{Constantinou:2017sej} 计算。
此外，与唯象估计的比较通常在 $\bar\mu{=}2$ GeV 进行。这就定义了转换因子表达式中选择的
$\MSb$ 重正化标度的值。另一种做法是选取 $\bar\mu{=}\bar\mu_0$，然后把标度演化到 2 GeV，
这需要用到算符的反常量纲：
\be
\gamma_{\cal O} =  -3 \frac{g^2\,C_f}{16\pi^2} \,.
\ee
注意，在单圈水平上反常量纲不依赖算符的 Dirac 结构，并且在 RI$'$ 与 $\MSb$ 方案中相同。
演化到 $\MSb$ 重正化标度 $\bar\mu{=}2$ GeV 可以经由中间的重正化群不变方案（RGI）实现，
如文献 Refs.~\cite{Constantinou:2014fka,Alexandrou:2015sea} 中所采用的。在本文的框架中，
我们对两种方法都做了检验，发现它们的差异可以忽略。

\begin{figure}[h]
\centering
\includegraphics[scale=0.265,angle=-90]{CV_CA_RIMS_Landau_pL3_Re.eps}\,\,\,
\includegraphics[scale=0.265,angle=-90]{CV_CA_RIMS_Landau_pL3_Im.eps}
\vspace*{-0.3cm}
\begin{minipage}{15cm}
\hspace*{3cm}
\caption{\small{$\bar\mu{=}\bar\mu_0$ 时从 RI$'$ 方案到 $\MSb$ 方案的单圈转换因子。
左（右）图显示转换因子的实部（虚部）作为以格点单位表示的 Wilson 线长度的函数。
采用了两种 RI$'$ 标度的选择：$a\bar\mu_0{=}\frac{2\pi}{32}\, (\frac{7}{2}{+}\frac{1}{4}, 3,3,3)$
（蓝色圆圈）与 $a\bar\mu_0{=}\frac{2\pi}{32}\, (\frac{4}{2}{+}\frac{1}{4}, 0,0,3)$（红色方块）。
在图例中分别使用缩写 (7,3,3,3) 与 (4,0,0,3)。}}
\label{fig:C}
\end{minipage}
\end{figure}

\medskip
如上一节所讨论的，关于空间方向我们采用两类 RI$'$ 重正化标度：一种与核子矩阵元中的相同，即
$a\vec{\bar{\mu}}_0{=}\frac{2\pi}{L}\,(0,0,P_3)$（平行于 Wilson 线的方向）；另一种是对角的、每个方向都取值 $P_3$，
即 $a\vec{\bar{\mu}}_0{=}\frac{2\pi}{L}\,(P_3,P_3,P_3)$。转换因子既依赖 RI$'$ 也依赖 $\MSb$ 重正化标度。
$\bar\mu$ 通常固定为 2 GeV，而 $\bar\mu_0$ 可以变化，从而影响转换因子的数值。图~\ref{fig:C} 展示了
非极化和螺旋度算符的这一情形，二者共享同一个转换因子\footnote{单圈计算不依赖于人们把 $\gamma_5$
推广到 $D$ 维时所采用的约定（四种相关约定及它们之间部分转换因子的讨论，参见例如
Refs.~\cite{Buras:1989xd,Patel:1992vu,Larin:1993tp,Larin:1993tq,Skouroupathis:2008mf,Constantinou:2013pba}）。}。
我们聚焦于核子以 $a P_3{=}\frac{6\pi}{L}$ 增强时的矩阵元重正化，并对转换因子在两种重正化标度情形下做同样的处理。
对该特定的 $P_3$ 值，我们选择 $\bar\mu_0$ 的时间分量使得上面定义的比值 $\hat{P}$ 尽可能小，
以压制格点伪影。尽管如此，对``平行''情形该比值的最小值为 $\hat{P}{=}0.54$，已经很高。
对``对角''和``平行''情形所选的 $a\bar\mu_0$ 值分别为
$\frac{2\pi}{32}\, (\frac{7}{2}{+}\frac{1}{4}, 3,3,3)$ 与
$\frac{2\pi}{32}\, (\frac{4}{2}{+}\frac{1}{4}, 0,0,3)$。

\medskip
转换因子的实部和虚部作为 Wilson 线长度 $z$ 的函数绘出。我们强调，转换因子对重正化标度和
Wilson 线长度的依赖高度非平凡，它用修正贝塞尔函数的积分来表达。因此，图~\ref{fig:C} 所示的数据点是
针对具体标度在每个 $z$ 值处分别数值计算的。我们观察到转换因子的实部比虚部大约一个数量级。
实部随 $z$ 的增加持续增大，而虚部在 $z$ 变为非零后几乎立即稳定下来。另外，$z{=}0$ 处的转换因子等于 1，
因为它对应局域矢量流和轴向流。注意 $z{=}0$ 的情形无法从文献~\cite{Constantinou:2017sej} 的计算中取得，
因为那里严格限于 $z{\neq}0$\,：超出树图水平的收缩项的出现使极限 $z\to 0$ 成为非解析的。相反，
上一节的非微扰处方可用于任意 $z$ 值，因为计算是对每个 $z$ 分别进行的。图~\ref{fig:C} 中的数值将在下一节中用于把
RI$'$ 非微扰 $Z$ 因子带到 $\MSb$ 方案。为了可靠地提取 $Z$ 因子，我们把计算扩展到若干个 $\bar\mu_0$ 值，
如第~\ref{Sec3} 节所解释。我们发现转换因子对所有 $\bar\mu_0$ 值都具有相同的定性行为。
